\documentclass[../DoAn.tex]{subfiles}
\begin{document}

\section{Thiết kế kiến trúc}

\subsection{Lựa chọn kiến trúc phần mềm}

Hệ thống trong đồ án là một ứng dụng web có nhiều nhóm chức năng khác nhau, bao gồm quản lý tenant, quản lý chatbot, quản lý knowledge base, xử lý hội thoại, tạo lead, tạo yêu cầu mua hàng, tích hợp kênh chat ngoài và xử lý các tác vụ AI. Các nhóm chức năng này có tính chất khác nhau: một phần thuộc nghiệp vụ quản trị và lưu trữ dữ liệu, một phần thuộc xử lý hội thoại và truy xuất tri thức, một phần thuộc giao diện người dùng. Vì vậy, hệ thống cần một kiến trúc có khả năng tách biệt các thành phần theo trách nhiệm, giúp dễ phát triển, kiểm thử và mở rộng.

Trong phạm vi đồ án, kiến trúc được lựa chọn là kiến trúc phân lớp kết hợp tách service theo chức năng. Cụ thể, hệ thống được chia thành các thành phần chính gồm frontend, backend nghiệp vụ, AI service và cơ sở dữ liệu. Frontend chịu trách nhiệm hiển thị giao diện và tiếp nhận thao tác của người dùng. Backend nghiệp vụ chịu trách nhiệm xử lý xác thực, phân quyền, quản lý tenant, chatbot, knowledge base, hội thoại, lead và yêu cầu mua hàng. AI service chịu trách nhiệm xử lý hội thoại, truy xuất knowledge base và sinh câu trả lời tư vấn. Cơ sở dữ liệu lưu trữ các dữ liệu nghiệp vụ chính của hệ thống.

Cách tổ chức này có điểm gần với kiến trúc ba lớp truyền thống, trong đó lớp giao diện tương ứng với frontend, lớp xử lý nghiệp vụ tương ứng với backend và AI service, còn lớp dữ liệu tương ứng với cơ sở dữ liệu và knowledge base. Tuy nhiên, do hệ thống có thành phần AI tương đối độc lập, phần xử lý AI được tách riêng thành một service thay vì đặt trực tiếp trong backend nghiệp vụ. Cách tách này giúp backend tập trung vào nghiệp vụ quản lý và phân quyền, còn AI service tập trung vào xử lý ngôn ngữ tự nhiên, truy xuất tri thức và sinh phản hồi.

Một số kiến trúc khác cũng có thể được cân nhắc. Kiến trúc nguyên khối có ưu điểm là đơn giản khi triển khai ban đầu, nhưng nếu đặt toàn bộ giao diện, nghiệp vụ, AI và lưu trữ trong một ứng dụng duy nhất thì hệ thống sẽ khó bảo trì khi số lượng chức năng tăng. Kiến trúc microservice đầy đủ cho phép tách rất nhiều service nhỏ, nhưng làm tăng độ phức tạp về triển khai, giao tiếp và giám sát, chưa thật sự cần thiết trong phạm vi đồ án. Kiến trúc MVC phù hợp để tổ chức một ứng dụng web đơn lẻ, nhưng chưa thể hiện rõ việc tách riêng AI service và pipeline RAG. Vì vậy, phương án phân lớp kết hợp tách AI service được lựa chọn vì cân bằng giữa tính rõ ràng, khả năng mở rộng và mức độ phức tạp phù hợp với đồ án tốt nghiệp.

Trong hệ thống cụ thể của đồ án, frontend cung cấp các màn hình đăng nhập, quản trị hệ thống, quản trị cửa hàng, quản lý chatbot, quản lý knowledge base, giao diện chat, danh sách hội thoại, lead và yêu cầu mua hàng. Backend cung cấp các API cho frontend và các webhook từ kênh chat ngoài. Backend cũng là nơi kiểm soát xác thực, phân quyền và tenant isolation trước khi truy cập dữ liệu. AI service nhận yêu cầu hội thoại từ backend, thực hiện truy xuất knowledge base tương ứng với tenant và trả về câu trả lời tư vấn. Cơ sở dữ liệu lưu thông tin tenant, người dùng, chatbot, hội thoại, tin nhắn, lead, yêu cầu mua hàng và trạng thái xử lý dữ liệu.

Với kiến trúc này, dữ liệu và cấu hình của từng cửa hàng được gắn với tenant tương ứng. Khi người dùng chat với chatbot của một cửa hàng, backend xác định tenant và chatbot liên quan, sau đó chuyển yêu cầu sang AI service cùng với thông tin ngữ cảnh cần thiết. AI service chỉ sử dụng knowledge base phù hợp với tenant đó để phục vụ quá trình tư vấn. Cách tổ chức này giúp hệ thống đáp ứng yêu cầu multi-tenant đã xác định ở Chương 2, đồng thời vẫn giữ được sự tách biệt giữa phần nghiệp vụ và phần xử lý AI.

\subsection{Thiết kế tổng quan}

Về tổng quan, hệ thống được tổ chức thành các nhóm thành phần chính gồm frontend, backend API, các module nghiệp vụ, AI service, tầng lưu trữ dữ liệu và các thành phần tích hợp bên ngoài. Các thành phần này được sắp xếp theo hướng phụ thuộc một chiều: frontend gửi yêu cầu đến backend; backend điều phối các module nghiệp vụ, truy cập cơ sở dữ liệu và gọi AI service khi cần; AI service xử lý truy xuất tri thức và sinh phản hồi; các kênh chat ngoài giao tiếp với hệ thống thông qua webhook do backend tiếp nhận.

Frontend là tầng giao diện của hệ thống. Thành phần này bao gồm các màn hình phục vụ người dùng cuối, nhân viên cửa hàng, quản trị cửa hàng và quản trị hệ thống. Các chức năng giao diện chính gồm đăng nhập, quản lý tenant, quản lý thành viên, quản lý chatbot, quản lý knowledge base, giao diện chat, quản lý hội thoại, quản lý lead, quản lý yêu cầu mua hàng và xem trạng thái vận hành. Frontend không truy cập trực tiếp cơ sở dữ liệu hoặc AI service, mà gửi yêu cầu đến backend thông qua API.

Backend API là điểm tiếp nhận các yêu cầu từ frontend và các kênh chat ngoài. Thành phần này chịu trách nhiệm định tuyến yêu cầu đến các module nghiệp vụ phù hợp, kiểm tra xác thực, kiểm tra quyền truy cập và chuẩn hóa dữ liệu đầu vào. Backend API phụ thuộc vào các module nghiệp vụ như quản lý tenant, quản lý chatbot, quản lý knowledge base, quản lý hội thoại, quản lý lead và quản lý yêu cầu mua hàng.

Thành phần xác thực và phân quyền phụ trách xác thực người dùng và phân quyền theo vai trò. Thành phần này giúp hệ thống phân biệt quản trị hệ thống, quản trị cửa hàng và nhân viên cửa hàng. Ngoài ra, thành phần này cũng hỗ trợ kiểm soát phạm vi tenant, bảo đảm người dùng thuộc tenant nào chỉ được truy cập dữ liệu thuộc tenant đó, trừ các chức năng dành riêng cho quản trị hệ thống.

Module quản lý tenant phụ trách thông tin cửa hàng và thành viên cửa hàng. Đây là module quan trọng đối với yêu cầu multi-tenant, vì các dữ liệu như chatbot, knowledge base, hội thoại, lead và yêu cầu mua hàng đều cần gắn với tenant tương ứng. Việc quản lý tenant rõ ràng giúp hệ thống tách biệt dữ liệu giữa các cửa hàng và hỗ trợ mở rộng thêm cửa hàng mới trong tương lai.

Module quản lý chatbot phụ trách tạo và cập nhật chatbot của từng cửa hàng. Các thông tin như tên chatbot, kênh sử dụng, persona, chế độ hoạt động và các cấu hình liên quan được quản lý trong module này. Mỗi chatbot thuộc về một tenant cụ thể, do đó mọi thao tác với chatbot đều cần được kiểm tra theo phạm vi tenant.

Module quản lý knowledge base quản lý các nguồn dữ liệu dùng để xây dựng knowledge base. Module này cho phép quản trị cửa hàng thêm nguồn dữ liệu, xóa nguồn dữ liệu, kích hoạt rebuild knowledge base và xem trạng thái xử lý. Đây là thành phần trung gian giữa dữ liệu của cửa hàng và quá trình tư vấn của chatbot, giúp chatbot có thể trả lời dựa trên dữ liệu riêng của từng tenant.

Module quản lý hội thoại lưu trữ và quản lý các phiên hội thoại giữa người dùng và chatbot. Module này lưu tin nhắn của người dùng, phản hồi của chatbot và các thông tin ngữ cảnh cần thiết. Dữ liệu hội thoại có vai trò quan trọng vì nó không chỉ phục vụ việc xem lại nội dung tư vấn, mà còn có thể được sử dụng để tạo lead hoặc yêu cầu mua hàng.

Module quản lý lead và yêu cầu mua hàng quản lý các thông tin phát sinh từ hội thoại tư vấn. Khi người dùng thể hiện ý định mua hàng hoặc cung cấp thông tin liên hệ, hệ thống có thể tạo lead hoặc yêu cầu mua hàng. Nhân viên cửa hàng sau đó có thể xem danh sách, xem chi tiết, claim, phân công và cập nhật trạng thái xử lý. Module này giúp kết nối quá trình tư vấn tự động với nghiệp vụ bán hàng của cửa hàng.

AI service phụ trách các chức năng xử lý ngôn ngữ tự nhiên. Thành phần này nhận yêu cầu từ backend, truy xuất dữ liệu liên quan trong knowledge base, sử dụng mô hình ngôn ngữ để sinh câu trả lời và trả kết quả về backend. Việc tách AI service thành một thành phần riêng giúp hệ thống dễ thay đổi mô hình, chiến lược truy xuất hoặc prompt mà không ảnh hưởng lớn đến phần nghiệp vụ.

Thành phần tích hợp kênh ngoài xử lý giao tiếp với các kênh như Messenger và Telegram. Thành phần này nhận webhook, chuẩn hóa tin nhắn về định dạng nội bộ, xác định tenant hoặc chatbot tương ứng và chuyển nội dung vào luồng xử lý hội thoại. Sau khi có câu trả lời, hệ thống gửi phản hồi lại kênh tương ứng.

Tầng lưu trữ dữ liệu làm việc với cơ sở dữ liệu để lưu tenant, người dùng, chatbot, nguồn dữ liệu knowledge base, hội thoại, tin nhắn, lead, yêu cầu mua hàng và các trạng thái vận hành. Các module nghiệp vụ không nên thao tác trực tiếp với cơ sở dữ liệu mà thông qua tầng truy cập dữ liệu để đảm bảo tính tổ chức, dễ kiểm soát và dễ thay đổi khi cần.

\subsection{Thiết kế chi tiết gói}

Trong phạm vi thiết kế chi tiết gói, đồ án tập trung vào ba nhóm package quan trọng nhất: nhóm quản lý chatbot và knowledge base, nhóm xử lý hội thoại tư vấn, và nhóm quản lý lead/yêu cầu mua hàng. Đây là các nhóm package gắn trực tiếp với các use case chính đã đặc tả ở Chương 2.

\subsubsection{Thiết kế gói Quản lý chatbot và knowledge base}

Gói quản lý chatbot và knowledge base chịu trách nhiệm cho các chức năng tạo chatbot, cập nhật cấu hình chatbot, quản lý nguồn dữ liệu knowledge base và rebuild knowledge base. Đây là nhóm chức năng được sử dụng bởi quản trị cửa hàng để chuẩn bị dữ liệu và cấu hình trước khi chatbot được đưa vào sử dụng.

Trong gói này, thành phần tiếp nhận yêu cầu có nhiệm vụ nhận các thao tác từ giao diện quản trị, chẳng hạn tạo chatbot mới, cập nhật cấu hình chatbot, thêm nguồn dữ liệu knowledge base hoặc kích hoạt rebuild. Các yêu cầu này được chuyển tới thành phần xử lý nghiệp vụ để kiểm tra dữ liệu đầu vào, kiểm tra quyền theo tenant và thực hiện thao tác tương ứng.

Thành phần xử lý nghiệp vụ chatbot chịu trách nhiệm quản lý các thông tin liên quan đến chatbot. Một chatbot cần được gắn với một tenant cụ thể để bảo đảm nó chỉ hoạt động trong phạm vi cửa hàng tương ứng. Các thông tin cấu hình của chatbot có thể bao gồm tên chatbot, kênh hoạt động, persona, chế độ hoạt động và các thông tin phục vụ việc kết nối với mô hình ngôn ngữ hoặc AI service nếu có.

Đối với knowledge base, thành phần xử lý nghiệp vụ cần quản lý danh sách nguồn dữ liệu của từng tenant. Khi quản trị cửa hàng thêm hoặc xóa nguồn dữ liệu, hệ thống cập nhật danh sách nguồn tương ứng. Khi quản trị cửa hàng yêu cầu rebuild, hệ thống ghi nhận yêu cầu xử lý, gọi thành phần AI service hoặc thành phần xử lý dữ liệu để xây dựng lại knowledge base, sau đó cập nhật trạng thái xử lý cho người dùng theo dõi.

Dữ liệu chính của gói này gồm chatbot, nguồn dữ liệu knowledge base và lịch sử rebuild. Chatbot biểu diễn cấu hình chatbot của một tenant. Nguồn dữ liệu knowledge base biểu diễn các tài nguyên được sử dụng để xây dựng tri thức cho chatbot. Lịch sử rebuild biểu diễn các lần xử lý dữ liệu, trạng thái thành công hoặc thất bại và thông tin lỗi nếu có. Cách tổ chức này giúp cửa hàng theo dõi được knowledge base đã sẵn sàng hay chưa trước khi đưa chatbot vào sử dụng.

\subsubsection{Thiết kế gói Xử lý hội thoại tư vấn}

Gói xử lý hội thoại tư vấn chịu trách nhiệm tiếp nhận tin nhắn từ người dùng, lưu hội thoại, gọi AI service để sinh câu trả lời và trả kết quả về giao diện hoặc kênh chat ngoài. Đây là gói trung tâm của hệ thống, vì nó trực tiếp phục vụ chức năng tư vấn sản phẩm theo cửa hàng.

Khi người dùng gửi tin nhắn, thành phần tiếp nhận yêu cầu chat xác định hội thoại hiện tại, chatbot tương ứng và tenant liên quan. Nếu người dùng bắt đầu một phiên mới, hệ thống tạo hội thoại mới. Nếu người dùng tiếp tục phiên cũ, hệ thống lấy lại thông tin hội thoại trước đó để duy trì ngữ cảnh. Tin nhắn của người dùng được lưu lại trước khi chuyển sang bước xử lý tư vấn.

Thành phần xử lý hội thoại chịu trách nhiệm điều phối toàn bộ quá trình tư vấn. Thành phần này lưu tin nhắn người dùng, gửi yêu cầu sang AI service, nhận câu trả lời từ AI service và lưu phản hồi của chatbot. Việc lưu cả tin nhắn người dùng và phản hồi của chatbot giúp hệ thống có lịch sử hội thoại đầy đủ, phục vụ việc xem lại, đánh giá chất lượng tư vấn và xử lý lead hoặc yêu cầu mua hàng.

AI service nhận thông tin hội thoại, tenant và chatbot từ backend. Dựa trên tenant tương ứng, AI service truy xuất dữ liệu liên quan từ knowledge base của cửa hàng, sau đó sử dụng kết quả truy xuất để sinh câu trả lời. Kết quả trả về có thể là câu trả lời tư vấn, gợi ý sản phẩm, câu hỏi làm rõ hoặc thông tin liên quan đến ý định mua hàng của người dùng.

Trong trường hợp người dùng thể hiện ý định mua hàng hoặc để lại thông tin liên hệ, gói xử lý hội thoại có thể phối hợp với gói quản lý lead hoặc gói quản lý yêu cầu mua hàng. Tuy nhiên, việc tạo và quản lý lead/yêu cầu mua hàng được tách sang gói nghiệp vụ riêng để tránh làm gói hội thoại trở nên quá lớn. Gói hội thoại chỉ đóng vai trò phát hiện bối cảnh và chuyển dữ liệu cần thiết sang module tương ứng.

Dữ liệu chính của gói này gồm hội thoại và tin nhắn. Hội thoại biểu diễn một phiên trao đổi giữa người dùng và chatbot. Tin nhắn biểu diễn từng lượt trao đổi trong hội thoại, bao gồm tin nhắn của người dùng và phản hồi của chatbot. Mỗi hội thoại cần gắn với tenant và chatbot để bảo đảm dữ liệu được quản lý đúng phạm vi.

\subsubsection{Thiết kế gói Quản lý lead và yêu cầu mua hàng}

Gói quản lý lead và yêu cầu mua hàng chịu trách nhiệm lưu trữ, hiển thị và cập nhật trạng thái các thông tin phát sinh từ hội thoại tư vấn. Gói này giúp kết nối chatbot với quy trình bán hàng của cửa hàng, thay vì để hội thoại chỉ dừng lại ở mức hỏi đáp.

Lead là thông tin khách hàng tiềm năng được tạo ra khi người dùng thể hiện nhu cầu mua hàng hoặc cung cấp thông tin liên hệ. Một lead cần lưu các thông tin như tenant, hội thoại liên quan, nhu cầu của khách hàng, sản phẩm quan tâm nếu có, thông tin liên hệ và trạng thái xử lý. Nhân viên cửa hàng có thể xem danh sách lead, xem chi tiết lead, xem lại hội thoại liên quan và cập nhật trạng thái xử lý.

Yêu cầu mua hàng là bản ghi thể hiện nhu cầu mua hàng rõ ràng hơn, cần được cửa hàng tiếp nhận và xử lý. Một yêu cầu mua hàng có thể được tạo trực tiếp từ hội thoại hoặc từ lead. Yêu cầu này cần lưu thông tin khách hàng, sản phẩm quan tâm, ghi chú, trạng thái xử lý và nhân viên phụ trách nếu đã được claim hoặc phân công.

Thành phần xử lý nghiệp vụ của gói này cần kiểm tra quyền truy cập theo tenant trước khi cho phép nhân viên xem hoặc cập nhật dữ liệu. Nhân viên cửa hàng chỉ được xem lead và yêu cầu mua hàng thuộc tenant của mình. Quản trị cửa hàng có thể có thêm quyền phân công lại yêu cầu mua hàng cho nhân viên khác trong cùng tenant.

Gói này có quan hệ chặt chẽ với gói xử lý hội thoại vì lead và yêu cầu mua hàng được sinh ra từ các phiên tư vấn. Khi nhân viên cần xử lý một lead hoặc yêu cầu mua hàng, hệ thống cần cho phép xem lại hội thoại liên quan để hiểu bối cảnh tư vấn trước đó. Gói này cũng có quan hệ với gói quản lý tenant và thành viên vì việc claim hoặc phân công yêu cầu mua hàng cần xác định nhân viên thuộc đúng cửa hàng.

Việc tách lead và yêu cầu mua hàng thành một gói riêng giúp hệ thống tổ chức rõ ràng các nghiệp vụ sau hội thoại. Chatbot chịu trách nhiệm tư vấn và ghi nhận nhu cầu, còn nhân viên cửa hàng chịu trách nhiệm xử lý các cơ hội bán hàng đã được hệ thống ghi nhận. Cách thiết kế này giúp hệ thống phù hợp hơn với quy trình vận hành thực tế của cửa hàng nội thất.

% \section{Thiết kế chi tiết}
% \subsection{Thiết kế giao diện}
% \subsection{Thiết kế lớp}
%
% \subsection{Thiết kế cơ sở dữ liệu}
% \section{Xây dựng ứng dụng}
% \subsection{Thư viện và công cụ sử dụng}
% \subsection{Kết quả đạt được}
% \subsection{Minh họa các chức năng chính}
% \section{Kiểm thử}
% \section{Triển khai}

\end{document}